% ============================================================
% VERSIÓN INDIVIDUAL — DANIEL DUARTE CORDERO
% Eje A completo con las lecturas asignadas a Daniel.
% Ejes B y C: se rellena únicamente lo sostenido por sus fuentes;
% los TODO indican material que debe integrarse con Adriel/Sebastián.
% ============================================================
\section{Trabajos Relacionados}
\label{sec:relacionados}

Para descubrimiento y verificación se consultaron
fuentes editoriales primarias y bases bibliográficas de IEEE, ACM,
AAAI/AIIDE, Springer y los proceedings correspondientes. Se priorizaron
trabajos arbitrados sobre predicción de victoria durante la partida,
evaluación de estados en juegos de cartas, generalización entre
composiciones o parches, y dos requisitos metodológicos del Track A:
imputación y desbalance. Se excluyeron estudios exclusivamente
pre-partida cuando no aportaban a la representación de estados ni a la
generalización. \textbf{\% TODO integración:} sustituir este párrafo por
la metodología de búsqueda del equipo completo, incluyendo cadenas de
búsqueda y el número total de trabajos cribados y retenidos.

\subsection{Eje A: el problema y el dominio}
\label{subsec:ejea}

La predicción de resultados en juegos competitivos aparece de forma
recurrente en la literatura de esports. La revisión sistemática de
Kamal et al. incluye 35 estudios de MOBA y encuentra un predominio del
aprendizaje supervisado, con uso frecuente de estadísticas de jugador,
mecánicas del juego y eventos intra-partida para tareas que incluyen la
predicción del ganador \cite{kamal_moba_2025}. El estudio de Tian et al.
es un ejemplo específico de predicción durante la partida: su Hero
Featured Network combina información de combate y de héroes a lo largo
del tiempo y reporta 81.8\% de accuracy en el nodo de 7 minutos en la
tabla detallada, aunque el resumen del artículo menciona 84.7\%
\cite{tian_hero_2022}. Más importante para nuestro problema, el estudio
muestra que la utilidad de las variables cambia con el progreso: la
alineación domina al inicio, el oro gana relevancia en la etapa media y
las bajas y torres se vuelven más informativas más adelante.

En juegos de cartas, la evaluación de estados se ha abordado con
objetivos distintos. Miernik y Kowalski evolucionan funciones lineales
y basadas en árboles para asignar un valor escalar a estados de
\emph{Legends of Code and Magic}; su mejor variante global alcanza
60.2\% de win rate en el torneo de evaluación, pero la función se
optimiza indirectamente por el desempeño del agente y no representa
una probabilidad de victoria aprendida de pares estado--resultado
\cite{miernik_evolving_2022}. El antecedente más cercano a nuestra tarea
es el AAIA'17 Data Mining Challenge de Hearthstone: su objetivo fue
predecir la probabilidad de victoria desde estados intra-partida y la
métrica del reto fue AUC. El baseline oficial alcanzó 0.7846 y la mejor
solución 0.8019 \cite{janusz_hearthstone_2017}. Este antecedente
establece que el problema es viable en otro CCG, pero no resuelve su
transferencia a Pokémon TCG ni sus riesgos particulares de dependencia
entre snapshots.

Otra familia de trabajos se concentra menos en la precisión dentro de
una distribución fija y más en la robustez frente a cambios del
entorno. Chitayat et al. proponen Clustered Character Representation
(CCR), que codifica héroes de Dota~2 mediante atributos y habilidades
de diseño en lugar de depender solo de IDs. Un modelo con CCR obtiene
AUC 0.85 en parches conocidos y mantiene 0.85 y 0.86 en dos parches
posteriores no usados para entrenar el modelo con partidas, lo que
apoya el uso de representaciones ligadas a mecánicas cuando el juego
cambia \cite{chitayat_meta_2023}. Sin embargo, su caso de estudio
predice conteos finales de bajas usando duración y composición de
equipo, no probabilidad de victoria desde un estado intermedio rico.

VGC-Bench aborda otra forma de cambio: la diversidad de composiciones.
El benchmark de Pokémon VGC separa equipos vistos y 72 equipos fuera de
distribución, con 1000 episodios por emparejamiento en sus pruebas
principales. En la prueba de generalización, el agente entrenado con 64
equipos vence al entrenado con uno con win rate 0.669, y los autores
concluyen que una mayor diversidad de entrenamiento mejora la
transferencia a equipos no vistos aun cuando empeora el rendimiento en
configuraciones familiares \cite{angliss_vgcbench_2026}. La tarea es de
aprendizaje de políticas, no de estimación probabilística, pero el
diseño de holdout por composición ofrece un antecedente metodológico
para preguntar si nuestro predictor aprende estado de juego o memoriza
el metajuego.

\noindent\textbf{Qué enfoques dominan.} En esports predominan métodos
supervisados para predicción de resultado y, en trabajos de estado
dinámico, representaciones que incorporan información intra-partida.
En CCG y Pokémon también aparecen evaluadores heurísticos, búsqueda,
behavior cloning y aprendizaje por refuerzo, pero esos enfoques
optimizan decisiones o políticas y no necesariamente producen
probabilidades calibradas \cite{kamal_moba_2025,miernik_evolving_2022,angliss_vgcbench_2026}.

\noindent\textbf{Qué se considera resuelto.} La literatura demuestra
que es posible extraer señal predictiva del estado durante la partida y
que la importancia de las variables puede cambiar con la progresión
temporal \cite{tian_hero_2022,janusz_hearthstone_2017}. También existen
protocolos explícitos para estudiar robustez a parches y a
composiciones no vistas \cite{chitayat_meta_2023,angliss_vgcbench_2026}.

\noindent\textbf{Qué sigue abierto.} No se ha establecido, con las
referencias revisadas, un protocolo equivalente para Pokémon TCG que
combine snapshots tabulares observables, control de dependencia entre
estados de una misma partida, calidad probabilística y evaluación bajo
cambio temporal y de mazos. Tampoco es válido asumir que los resultados
de MOBA, VGC o Hearthstone se transfieren directamente porque cambian
las reglas, las unidades de observación, los protocolos y la población.

\subsection{Eje B: las técnicas del track}
\label{subsec:ejeb}

Dentro de las referencias asignadas a Daniel, dos trabajos sostienen
decisiones transversales del Track A. Emmanuel et al. revisan
mecanismos de datos faltantes y métodos de imputación, y muestran que
la precisión relativa de KNN y missForest cambia entre datasets
\cite{emmanuel_missing_2021}. El resultado relevante para este proyecto
no es que una técnica sea superior, sino que la estrategia de
imputación depende del mecanismo de ausencia y de las características
del conjunto de datos. Por ello, una frecuencia alta de valores
ausentes no justifica por sí sola imputarlos automáticamente.

Van den Goorbergh et al. estudian modelos de regresión logística
estándar y ridge bajo ausencia de corrección, random undersampling,
random oversampling y SMOTE. Las correcciones no mejoran de forma
sistemática el AUROC y sí producen sobreestimación pronunciada de la
probabilidad de la clase minoritaria, deteriorando la calibración
\cite{vandenGoorbergh_imbalance_2022}. El alcance debe respetarse: el
trabajo no evalúa \texttt{class\_weight}, Random Forest, SVM o redes
neuronales; por tanto, no justifica preferir ponderación de clases como
regla general.

\noindent\textbf{Fortalezas.} Estas referencias permiten justificar un
tratamiento de faltantes guiado por mecanismo y una evaluación del
desbalance centrada no solo en clasificación sino también en
calibración cuando la salida deseada es una probabilidad
\cite{emmanuel_missing_2021,vandenGoorbergh_imbalance_2022}.

\noindent\textbf{Límites.} Ninguna de estas dos referencias decide por
sí sola qué modelo tabular clásico debe ser el principal. Paper 6 no
formaliza nuestra noción de ausencia estructural/semántica, y Paper 7
solo estudia correcciones de desbalance en regresión logística.
\textbf{\% TODO integración:} incorporar las referencias de Adriel sobre
modelos tabulares, calibración y explicabilidad para completar las
3--5 síntesis exigidas por el Eje B.

\noindent\textbf{Supuestos bajo los que funcionan.} Las decisiones de
imputación requieren identificar razonablemente el mecanismo y tipo de
dato; las conclusiones de Paper 7 son directamente aplicables a modelos
de regresión logística probabilística y solo indicativas para otras
familias.

\noindent\textbf{Justificación técnica del track.} Con las referencias
de Daniel puede sostenerse que un problema de probabilidad sobre
características estructuradas exige especial cuidado con calidad de
datos y calibración. La justificación completa de por qué se elige
Track A frente a una alternativa profunda o de RL debe integrarse con
las referencias metodológicas del resto del equipo. \textbf{\% TODO
integración.}

\subsection{Eje C: el dataset}
\label{subsec:ejec}

El dataset \emph{PTCG AI Battle Challenge Simulation Episodes} fue
publicado en 2026 y, en la revisión realizada hasta esta etapa, no se
localizó literatura arbitrada que lo utilice. Por ello, el requisito de
baseline se sustituye por datasets y entornos análogos. Hearthstone es
el análogo más cercano porque también es un juego de cartas con
información imperfecta y el reto AAIA'17 predice el ganador desde un
estado intra-partida \cite{janusz_hearthstone_2017}. LoCM aporta un
antecedente conceptual de evaluación de estados en CCG, aunque su
protocolo es simulación evolutiva y su salida no es probabilística
\cite{miernik_evolving_2022}.

VGC-Bench y el trabajo de Chitayat et al. no son datasets equivalentes,
pero aportan protocolos útiles para dos riesgos presentes en nuestros
datos: composiciones no vistas y cambio de versión. VGC-Bench reserva
configuraciones de equipo fuera de entrenamiento, mientras que CCR se
evalúa en dos parches posteriores excluidos del entrenamiento con
partidas \cite{angliss_vgcbench_2026,chitayat_meta_2023}. Estas
comparaciones justifican incorporar pruebas de robustez separadas del
holdout principal.

\noindent\textbf{Métricas y protocolos reportados.} Hearthstone usa AUC
para comparar predictores de probabilidad; Tian et al. reportan
accuracy por nodos temporales; Chitayat et al. reportan AUC entre
parches; VGC-Bench usa win rate, cross-play, Alpha-Rank y evaluaciones
sobre equipos vistos/no vistos \cite{janusz_hearthstone_2017,tian_hero_2022,chitayat_meta_2023,angliss_vgcbench_2026}. La heterogeneidad impide
leer sus números como un ranking único.

\noindent\textbf{Problemas conocidos.} Los análogos revelan riesgos de
representatividad y generalización: Tian et al. usan un solo día y
jugadores de alto nivel; Chitayat et al. restringen sus partidas a
profesionales/premium; VGC-Bench muestra degradación al ampliar la
diversidad de equipos. En el dataset propio, además, la auditoría
muestra fuerte dependencia entre snapshots de un mismo episodio y
cambios temporales de mazos, agentes y versiones. \textbf{\% TODO
integración:} incorporar aquí la revisión específica de Sebastián sobre
el dataset y análogos adicionales.

\subsection{Análisis comparativo}
\label{subsec:comparativo}

\safeinput{tab/tabla_comparativa}

La tabla muestra que los resultados numéricos no admiten una
comparación directa. Tian et al. evalúan accuracy de un modelo
secuencial en un MOBA, Miernik y Kowalski evalúan win rate de agentes
que consumen una función heurística, VGC-Bench compara políticas por
cross-play y el estudio de Chitayat et al. usa AUC para una tarea de
predicción de bajas. Hearthstone es el único antecedente cuyo objetivo
y métrica se aproximan a nuestro planteamiento: estimar la posibilidad
de victoria desde un estado intra-partida mediante AUC
\cite{janusz_hearthstone_2017}.

Sí aparece una tendencia transversal: el desempeño depende del
momento del juego y de la distribución sobre la que se evalúa. Tian et
al. encuentran distinta utilidad de variables por etapa;
VGC-Bench separa composiciones vistas de no vistas; Chitayat et al.
separan parches de entrenamiento y parches futuros
\cite{tian_hero_2022,angliss_vgcbench_2026,chitayat_meta_2023}. Por ello,
un único holdout aleatorio sería insuficiente para sostener una
conclusión fuerte de generalización en Pokémon TCG.

Otra limitación recurrente es que el resultado depende de una
población acotada. Los estudios de HoK y Dota~2 revisados se restringen
a jugadores de alto nivel/profesionales, y VGC-Bench trabaja con un
entorno y formato distintos de Pokémon. Esto obliga a declarar los
límites de extrapolación y evita presentar un resultado sobre agentes
artificiales de Pokémon TCG como evidencia sobre jugadores humanos.

\subsection{Brechas identificadas}
\label{subsec:brechas}

\textbf{Carencias encontradas en la literatura.}
\begin{enumerate}
  \item No se localizó literatura arbitrada sobre el dataset específico
  de replays de la competencia Pokémon TCG de 2026.
  \item La estimación de victoria desde estados existe en Hearthstone,
  pero no está establecida para la estructura de recursos y estados de
  Pokémon TCG en las referencias revisadas.
  \item El antecedente de Hearthstone usa AUC como medida de
  discriminación, pero la calidad y calibración de la probabilidad
  requieren evaluación adicional cuando la salida se consume como
  evaluador de posiciones.
  \item Los datasets con múltiples estados correlacionados por partida
  requieren particiones por unidad independiente; en el conjunto propio
  un split por snapshot produciría fuga entre episodios.
  \item La literatura revisada muestra sensibilidad a composiciones y
  parches, pero no existe un protocolo único que combine partida nueva,
  mazo no visto y periodo futuro para esta tarea de Pokémon TCG.
\end{enumerate}

\noindent\textbf{Brechas que este trabajo aborda.} El proyecto aborda
las brechas 2--5 mediante una formulación supervisada de
$P(\mathrm{win}\mid\mathrm{state})$, métricas de discriminación y
calidad probabilística, particiones agrupadas por episodio, análisis
por turno y pruebas secundarias temporales y por mazo/matchup.

\noindent\textbf{Brechas fuera del alcance.} La ausencia de literatura
sobre el dataset no puede resolverse dentro del semestre; se documenta
y se sustituye por análogos. Tampoco se busca demostrar transferencia a
jugadores humanos, resolver el juego ni construir una política de RL.
La comparabilidad estricta con Hearthstone permanece fuera del alcance
porque cambian reglas, población, representación y protocolo.
